\chapter{1977年福建卷（文）}

\section{解答题}

\begin{problem}
\begin{enumerate}
    \item 计算：\( 5-3\times\left[ \left( -3\frac{3}{8} \right)^{-\frac{1}{3}}+1031\times\left( 0.25-2^{-2} \right) \right]\div 9^0 \).

    \item 求 \( \cos\left( -840^\circ \right) \) 的值.

    \item 化简：\( \sqrt{\left( 2x-3 \right)^2} \).

    \item 如图，在 \( \triangle ABC\) 中，\( MN\parallel BC\)，\( MN=1\mathrm{~cm}\)，\( BC=3\mathrm{~cm}\)，\( BM=AM+2\)，求 \(AM\) 的长.

    \item 已知 \( \lg 3=0.4771 \)，\( \lg x=3.4771 \)，求 \( x \).

    \item 求 \( \displaystyle\lim_{x\to 1}\frac{x-1}{x^2-3x+2}\).

    \item 求函数 \( y=x^2+2x-4 \) 的极小值.

    \item 已知 \( \sin \alpha=\frac{3}{5} \)，\( \frac{\pi}{2}<\alpha<\pi \)，求 \( \tan \alpha \) 的值.

    \item 写出等比数列 \( -\frac{2}{9} \)，\( \frac{2}{27} \)，\( -\frac{2}{81} \)，\( \cdots \) 的通项公式.
\end{enumerate}

\bitmapfigure[width=6.0cm]{img/1977/fujian_liberal/q1_fig1.png}

\end{problem}

\begin{solution}
\begin{enumerate}
    \item 因为 \( -3\frac{3}{8}=-\frac{27}{8}=\left(-\frac{3}{2}\right)^3 \)，所以 \( \left(-3\frac{3}{8}\right)^{-\frac{1}{3}}=-\frac{2}{3}\). 又因为 \( 0.25=\frac{1}{4}=2^{-2}\)，所以 \( 1031\left(0.25-2^{-2}\right)=0\)，且 \( 9^0=1\). 因此原式
    \[
    =5-3\left(-\frac{2}{3}\right)\div1=7
    \]

    \item \( \cos\left(-840^\circ\right)=\cos840^\circ=\cos\left(2\times360^\circ+120^\circ\right)=\cos120^\circ=-\frac{1}{2}\).

    \item 由算术平方根的定义，\( \sqrt{\left(2x-3\right)^2}=\left|2x-3\right|\). 当 \( 2x-3\geq0\)，即 \(x\geq\frac{3}{2}\) 时，原式 \( =2x-3\)；当 \( 2x-3<0\)，即 \(x<\frac{3}{2}\) 时，原式 \( =-(2x-3)=3-2x\)，故
    \[
    \sqrt{\left(2x-3\right)^2}=\begin{cases}
    2x-3,&x\geq\frac{3}{2},\\
    3-2x,&x<\frac{3}{2}.
    \end{cases}
    \]

    \item 因为 \( MN\parallel BC\)，所以 \( \triangle AMN\sim\triangle ABC\)，从而
    \[
    \frac{AM}{AB}=\frac{MN}{BC}=\frac{1}{3}
    \]
    又 \( AB=AM+BM=AM+(AM+2)=2AM+2\)，故
    \[
    \frac{AM}{2AM+2}=\frac{1}{3}\Longrightarrow 3AM=2AM+2\Longrightarrow AM=2\mathrm{~cm}
    \]

    \item 因为 \( 3.4771=0.4771+3=\lg3+\lg1000=\lg3000\)，所以 \( \lg x=\lg3000\). 对数函数在其定义域上是单调函数，故 \(x=3000\).

    \item 因为 \(x^2-3x+2=(x-1)(x-2)\)，当 \(x\ne1\) 时，
    \[
    \frac{x-1}{x^2-3x+2}=\frac{x-1}{(x-1)(x-2)}=\frac{1}{x-2}
    \]
    因而 \( \displaystyle\lim_{x\to1}\frac{x-1}{x^2-3x+2}=\lim_{x\to1}\frac{1}{x-2}=-1\).

    \item 配方得 \(y=x^2+2x-4=(x+1)^2-5\). 因为 \((x+1)^2\geq0\)，所以 \(y\geq-5\)，当且仅当 \(x=-1\) 时等号成立，故函数的极小值为 \(-5\).

    \item 因为 \( \frac{\pi}{2}<\alpha<\pi\)，所以 \( \alpha\) 在第二象限，\( \cos\alpha<0\). 由 \( \sin^2\alpha+\cos^2\alpha=1\)，得
    \[
    \cos\alpha=-\sqrt{1-\left(\frac{3}{5}\right)^2}=-\frac{4}{5}
    \]
    因此 \( \tan\alpha=\frac{\sin\alpha}{\cos\alpha}=\frac{3/5}{-4/5}=-\frac{3}{4}\).

    \item 该等比数列的首项为 \(a_1=-\frac{2}{9}\)，公比为
    \[
    q=\frac{\frac{2}{27}}{-\frac{2}{9}}=-\frac{1}{3}
    \]
    所以其通项公式为 \(a_n=a_1q^{n-1}=-\frac{2}{9}\left(-\frac{1}{3}\right)^{n-1}\)，其中 \(n=1\)，\(2\)，\(\ldots\).
\end{enumerate}
\end{solution}

\begin{problem}
\begin{enumerate}
    \item 求函数 \( y=\frac{\lg\left( 2-x \right)}{\sqrt{x-1}} \) 的定义域.

    \item 证明：\( \left( \sin \alpha-\cos \alpha \right)^2+\sin 2\alpha=1 \).

    \item 解方程：\( 2\sqrt{x-3}+6=x \).

    \item 解不等式：\( x^2-x-6<0 \).

    \item 把分母有理化：\( \sqrt{\frac{\sqrt{5}+\sqrt{2}}{\sqrt{5}-\sqrt{2}}} \).

    \item 某中学革命师生自己动手油漆一个直径为 \( 1.2\text{ 米}\) 的地球仪，如果每平方米面积需要油漆 \( 150\text{ 克}\)，问共需油漆多少克？（答案保留整数）
\end{enumerate}
\end{problem}

\begin{solution}
\begin{enumerate}
    \item 由对数的真数必须为正，得 \( 2-x>0\)，即 \(x<2\). 又分母中的根式必须有意义且不能为零，得 \(x-1>0\)，即 \(x>1\). 两条件取交集，函数的定义域为 \(\left(1,2\right)\).

    \item \( \left(\sin\alpha-\cos\alpha\right)^2+\sin2\alpha=\sin^2\alpha-2\sin\alpha\cos\alpha+\cos^2\alpha+2\sin\alpha\cos\alpha=\sin^2\alpha+\cos^2\alpha=1\). 得证.

    \item 由根式有意义，得 \(x\geq3\). 令 \(t=\sqrt{x-3}\)，则 \(t\geq0\)，且 \(x=t^2+3\). 代入原方程得
    \[
    2t+6=t^2+3\Longrightarrow t^2-2t-3=0\Longrightarrow \left(t-3\right)\left(t+1\right)=0
    \]
    因为 \(t\geq0\)，所以只能取 \(t=3\)，从而 \(x=t^2+3=12\). 代回原方程，\( 2\sqrt{12-3}+6=2\times3+6=12\)，故原方程的解为 \(x=12\).

    \item 因为 \(x^2-x-6=\left(x-3\right)\left(x+2\right)\)，不等式化为 \(\left(x-3\right)\left(x+2\right)<0\). 两个根为 \(-2\) 和 \(3\)，二次项系数为正，所以乘积小于零时 \(x\) 在两根之间，故不等式的解集为 \(\left(-2,3\right)\).

    \item 因为 \( \sqrt5-\sqrt2>0\)，所以
    \[
    \sqrt{\frac{\sqrt5+\sqrt2}{\sqrt5-\sqrt2}}=\sqrt{\frac{\left(\sqrt5+\sqrt2\right)^2}{5-2}}=\frac{\sqrt5+\sqrt2}{\sqrt3}=\frac{\sqrt{15}+\sqrt6}{3}
    \]

    \item 地球仪的半径为 \(r=\frac{1.2}{2}=0.6\text{ 米}\)，表面积为 \( S=4\pi r^2=4\pi\times0.6^2=1.44\pi\text{ 平方米}\). 所需油漆质量为
    \[
    150\times1.44\pi=216\pi\text{ 克}\approx678.6\text{ 克}
    \]
    按题意保留整数，得共需油漆 \( 679\text{ 克}\).
\end{enumerate}
\end{solution}

\begin{problem}
某农机厂开展“工业学大庆”运动，在十月份生产拖拉机 \( 1000\text{ 台}\). 这样，一月至十月的产量恰好完成全年生产任务. 工人同志为了加速农业机械化，计划在年底前再生产 \( 2310\text{ 台}\).

\begin{enumerate}
    \item 求十一月、十二月份每月增长率；
    \item 原计划年产拖拉机多少台？
\end{enumerate}
\end{problem}

\begin{solution}
\begin{enumerate}
    \item 设每月增长率为 \(x\)，则十一月份和十二月份的产量分别为 \( 1000\left(1+x\right)\text{ 台}\) 和 \( 1000\left(1+x\right)^2\text{ 台}\). 由计划在年底前再生产 \( 2310\text{ 台}\)，得
    \[
    1000\left(1+x\right)+1000\left(1+x\right)^2=2310
    \]
    令 \(u=1+x\)，则 \(u>0\)，且
    \[
    u+u^2=2.31\Longrightarrow u^2+u-2.31=0\Longrightarrow \left(u-1.1\right)\left(u+2.1\right)=0
    \]
    因为 \(u>0\)，所以 \(u=1.1\)，从而 \(x=0.1=10\%\).

    \item 第（1）问得到十一月份、十二月份的产量分别为 \(1000\times1.1=1100\text{ 台}\) 和 \(1000\times1.1^2=1210\text{ 台}\). 但是，题面只说明十月份产量为 \(1000\text{ 台}\)，以及一月至十月的总产量等于原计划年产量，没有说明一月至十月各月产量之间的关系. 因此设一月至十月的产量为 \(a_1\)，\(a_2\)，\(\ldots\)，\(a_{10}\)，已知的只有 \(a_{10}=1000\)，而原计划年产量为 \(a_1+a_2+\cdots+a_{10}\)，其中 \(a_1\)，\(\ldots\)，\(a_9\) 均没有被题面确定. 例如，取前九个月产量均为 \(100\text{ 台}\)，原计划年产量为 \(9\times100+1000=1900\text{ 台}\)；取前九个月产量均为 \(200\text{ 台}\)，原计划年产量为 \(9\times200+1000=2800\text{ 台}\). 这两种情形都满足题面条件，且不影响十一月、十二月份按第（1）问所得的生产计划，所以原计划年产拖拉机的台数不能由现有题面条件唯一确定.
\end{enumerate}
\end{solution}

\begin{problem}
求抛物线 \( y^2=9x \) 和圆 \( x^2+y^2=36 \) 在第一象限的交点处的切线方程.
\end{problem}

\begin{solution}
由 \(y^2=9x\)，代入圆的方程得 \(x^2+9x=36\)，即
\[
    \left(x+12\right)\left(x-3\right)=0
\]
交点在第一象限，故 \(x=3\)，并且 \(y^2=27\)，所以交点为 \( \left(3,3\sqrt3\right)\).

对抛物线方程两边关于 \(x\) 求导，得 \( 2yy'=9\)，所以在该点处切线斜率为
\[
y'=\frac{9}{2y}=\frac{9}{6\sqrt3}=\frac{\sqrt3}{2}
\]
因此切线方程为
\[
    y-3\sqrt3=\frac{\sqrt3}{2}\left(x-3\right)\Longrightarrow 3x-2\sqrt3y+9=0
\]
\end{solution}

\begin{problem}
已知双曲线 \( \frac{x^2}{24\alpha}-\frac{y^2}{16\cot\alpha}=1\)（\( \alpha\) 为锐角）和圆 \( \left(x-m\right)^2+y^2=r^2\) 相切于点 \( A\left(4\sqrt3,4\right)\)，求 \( \alpha\)，\(m\)，\(r\) 的值.
\end{problem}

\begin{solution}
将点 \( A\left(4\sqrt3,4\right)\) 代入双曲线方程，得
\[
\frac{48}{24\alpha}-\frac{16}{16\cot\alpha}=1\Longrightarrow \frac{2}{\alpha}-\tan\alpha=1
\]
令 \(F(t)=\frac{2}{t}-\tan t-1\). 因为 \(0<t<\frac{\pi}{2}\)，所以
\[
F'(t)=-\frac{2}{t^2}-\sec^2t<0
\]
又 \(\displaystyle\lim_{t\to0^+}F(t)=+\infty\)，\(\displaystyle\lim_{t\to(\pi/2)^-}F(t)=-\infty\)，故方程 \(F(t)=0\) 在 \(\left(0,\frac{\pi}{2}\right)\) 内有且只有一个根. 用二分法逐步缩小区间，例如
\(F(0.8924934185)\approx3.81\times10^{-11}>0\)，
\(F(0.89249341851)\approx-1.24\times10^{-11}<0\). 由于 \(F\) 严格递减，故
\(0.8924934185<\alpha<0.89249341851\)，从而 \(\alpha\approx0.89249342\)（弧度），代入正切函数得 \(\tan\alpha\approx1.24091288\).

对双曲线方程两边关于 \(x\) 求导，得
\[
\frac{x}{12\alpha}-\frac{yy'}{8\cot\alpha}=0
\]
所以在点 \(A\) 处的切线斜率为
\[
k=y'=\frac{2x\cot\alpha}{3\alpha y}=\frac{2\sqrt3\cot\alpha}{3\alpha}
\]
圆心为 \(C\left(m,0\right)\)，圆在点 \(A\) 处的半径垂直于切线，故半径 \(CA\) 的斜率为 \(-1/k\). 因而
\[
\frac{4}{4\sqrt3-m}=-\frac{1}{k}=-\frac{\sqrt3\alpha\tan\alpha}{2}
\]
由此 \(4k=m-4\sqrt3\)，所以
\[
m=4\sqrt3+4k=4\sqrt3+\frac{8}{\sqrt3\alpha\tan\alpha}\approx11.09865443
\]
又因为 \(4\sqrt3-m=-4k\)，故
\[
r^2=CA^2=\left(4\sqrt3-m\right)^2+4^2
=16k^2+16
=16\left(\frac{2\sqrt3}{3\alpha\tan\alpha}\right)^2+16
=\frac{64}{3\alpha^2\tan^2\alpha}+16
\]
所以
\[
r=\sqrt{16+\frac{64}{3\alpha^2\tan^2\alpha}}\approx5.77863853
\]
\end{solution}

\begin{problem}
某大队在农田基本建设的规划中，要测定被障碍物隔开的两点 \(A\) 和 \(P\) 之间的距离，他们土法上马，在障碍物的两侧，选取两点 \(B\) 和 \(C\)（如图），测得 \( AB=AC=50\mathrm{~m}\)，\( \angle BAC=60^\circ\)，\( \angle ABP=120^\circ\)，\( \angle ACP=135^\circ\)，求 \(A\) 和 \(P\) 之间的距离. （答案可用最简根式表示）

\bitmapfigure[width=6.4cm]{img/1977/fujian_liberal/q6_fig1.png}
\end{problem}

\begin{solution}
设 \( \theta=\angle BAP\). 由于 \( \angle BAC=60^\circ\)，有 \( \angle CAP=60^\circ-\theta\). 在 \( \triangle ABP\) 中，
\[
\angle APB=180^\circ-120^\circ-\theta=60^\circ-\theta
\]
在 \( \triangle ACP\) 中，
\[
\angle APC=180^\circ-135^\circ-\left(60^\circ-\theta\right)=\theta-15^\circ
\]
由正弦定理，在两个三角形中分别有
\(\frac{AP}{\sin120^\circ}=\frac{50}{\sin\left(60^\circ-\theta\right)}\)，\(\frac{AP}{\sin135^\circ}=\frac{50}{\sin\left(\theta-15^\circ\right)}\).
所以
\[
\frac{25\sqrt3}{\sin\left(60^\circ-\theta\right)}=\frac{25\sqrt2}{\sin\left(\theta-15^\circ\right)}
\Longrightarrow \sqrt3\sin\left(\theta-15^\circ\right)=\sqrt2\sin\left(60^\circ-\theta\right)
\]
由和差角公式，
\[
\sqrt3\left(\sin\theta\cos15^\circ-\cos\theta\sin15^\circ\right)
=\sqrt2\left(\sin60^\circ\cos\theta-\cos60^\circ\sin\theta\right)
\]
代入 \(\cos15^\circ=\frac{\sqrt6+\sqrt2}{4}\)，\(\sin15^\circ=\frac{\sqrt6-\sqrt2}{4}\)，\(\sin60^\circ=\frac{\sqrt3}{2}\)，\(\cos60^\circ=\frac{1}{2}\)，得
\[
    \frac{\sqrt3\left(\sqrt6+\sqrt2\right)}{4}\sin\theta-\frac{\sqrt3\left(\sqrt6-\sqrt2\right)}{4}\cos\theta
=\frac{\sqrt6}{2}\cos\theta-\frac{\sqrt2}{2}\sin\theta
\]
移项并合并同类项，利用 \(\sqrt3\sqrt6=3\sqrt2\)，得
\[
\frac{5\sqrt2+\sqrt6}{4}\sin\theta=\frac{3\sqrt2+\sqrt6}{4}\cos\theta
\]
因为 \(15^\circ<\theta<60^\circ\)，所以 \(\cos\theta>0\)，从而
\[
\tan\theta=\frac{3\sqrt2+\sqrt6}{5\sqrt2+\sqrt6}
=\frac{3+\sqrt3}{5+\sqrt3}
=\frac{6+\sqrt3}{11}
\]
再由
\[
\sin\left(60^\circ-\theta\right)=\frac{\sqrt3\cos\theta-\sin\theta}{2}
\]
并令 \(t=\tan\theta=\frac{6+\sqrt3}{11}\)，则
\(\cos^2\theta=\frac{1}{1+t^2}\)，故
\[
\sin^2\left(60^\circ-\theta\right)
=\frac{\left(\sqrt3-t\right)^2}{4\left(1+t^2\right)}
=\frac{\left(\sqrt3-\frac{6+\sqrt3}{11}\right)^2}{4\left(1+\left(\frac{6+\sqrt3}{11}\right)^2\right)}
=\frac{3}{10+4\sqrt3}
\]
因此
\[
    AP^2=\frac{\left(25\sqrt3\right)^2}{\sin^2\left(60^\circ-\theta\right)}=625\left(10+4\sqrt3\right)
\]
从而 \( AP=25\sqrt{10+4\sqrt3}\mathrm{~m}\).
\end{solution}
